In this section, we describe ways in which surfaces can be created in Orbiter. 

\bigskip

Orbiter contains a built-in catalogue of cubic surfaces with 27 
lines for small finite fields $\bbF_q.$ 
The surfaces in the catalogue all come with their automorphism group.
It is also possible to create surfaces from known families, 
or from associated objects like 6-arcs. 
Some of these constructions only create the surface, 
not the automorphism group.

\bigskip

Tables~\ref{tab:createsurface:one}-\ref{tab:createsurface:three} 
summarize the Orbiter commands that can be used
to create cubic surfaces. 
\orbitertableofcommands{surface_create_1}
\orbitertableofcommands{surface_create_2}
\orbitertableofcommands{surface_create_3}
The commands require a projective space object.
Not all of the surfaces created have 27 lines, 
and some of the constructions may yield degenerate surfaces. 




Let us look at some examples. 
The next command creates the unique surface with 27 lines 
over the field $\bbF_4.$ 
It is the Hirschfeld surface. 
In this command, 
the surface is pulled from Orbiter's built-in catalogue of cubic surfaces. 
The surface has Orbiter Catalogue Number (OCN) equal to 0.

\sourcecode{surface_4_0}


\bigskip


It is possible to create the same surface directly from the equation. 
At first, we define a variable holding the equation in algebraic form:


\sourcecodevariable{HIRSCHFELD_SURFACE_EQN_AF}

The following command creates the surface from the algebraic form and produces a report. 
In order to do so, the command creates a suitable polynomial ring $R$ and a symbolic object for the equation. 
Only then the surface is created and a report is written. 

\sourcecode{Hirschfeld_surface_q4}


\bigskip


There is a 4-parameter normal form $\cF_{abcd}$ for 
cubic surfaces with 27 lines. Here, $a,b,c,d$ are field elements in $\bbF_q$.
The following command creates the surface $\cF_{abcd}$ 
with $(a,b,c,d)=(2,3,3,4)$ over the field $\bbF_7$:
\sourcecode{Family_general_F7}
This command uses a built-in 
command to realize the surface of type $\cF_{abcd}$. 
We will see later that it is possible to 
create surfaces symbolically from an equation.

\bigskip


Another way of creating a cubic surface is as member of a known infinite family.
However, different choices of the parameters may lead to isomorphic surfaces.
For instance, 

\sourcecode{surface_eckardt_13_4_12}

creates the member of the Eckardt family  
described in~\cite{BettenKaraoglu19a} 
with parameters $(a,b) = (4,12)$ over the field $\bbF_{13}.$
On the other hand, the command 
\sourcecode{Eckardt_13}
creates the Eckardt surface over the field $\bbF_{13}$ 
from the parameter set $(a,b) = (3,1).$
The two surfaces are isomorphic.


\bigskip



It is possible to create a random cubic surface of order $q$. The following example 
creates a random cubic surface over the field $\bbF_{11}$:

\sourcecode{surface_11_random}

The only restriction is that $q$ is within the range 
for which the classification of cubic surfaces is available in Orbiter.

\bigskip

The normal form of~\cite{BettenKaraoglu22} 
can be used to create any cubic surface over $\bbF_q$.
This means that up to isomorphism, 
any cubic surface with 27 lines 
can be written as $\cF_{a,b,c,d}$ 
for a suitable choice of parameters $a,b,c,d \in \bbF_q.$
For instance, the following command uses $(a,b,c,d) = (2,30,30,2)$ 
to create the Eckardt surface over $\bbF_{31}$:


\sourcecode{surface_F_abcd_Eckardt_q31_by_equation}


\bigskip



A cubic surface is determined uniquely up 
to isomorphism by a non-conical 6-arc in a plane. 
The points of the arc are given in Orbiter ranks.
The following command creates the cubic surface which is 
associated with the arc $0,1,2,3,43,113$ over the field 
$\bbF_{13}$.


\sourcecode{Eckardt_13_by_arc_lifting}

The surface is of Eckardt type.

\bigskip


The formula for $\cF_{a,b,c,d}$ can be used symbolically. 
To this end, we encode it as a makefile varaible:

\sourcecodevariable{F_abcd_eqn}

We use notation such as $1+1+\cdots+1$ ($k$ times) 
for the integer $k$. This makes the formula valid 
over very small fields such as $\bbF_3$ and $\bbF_4,$
where the Orbiter conventions for field elements prohibit using $3$ or $4.$
There is no automatic reduction modulo 
the characteristic when specifying field elements.


\bigskip



To do so, we first create a symbolic object for the equation. 
Then we perform the substitution of $(a,b,c,d)$ with the desired values.
Finally, we create the surface from the equation:

\sourcecode{surface_F_abcd_Eckardt_q31_by_symbolic_object}



\bigskip

Using the Clebsch map, we can create a cubic surface from a non-conical 6-arc. 
The command expects the 6-arc in a plane in $\PG(3,q).$ 
The command also expects the rank numbers of two skew lines not in the plane contaoining the arc. 
If two rank numbers are given as two zeros, then Orbiter will pick a suitable set of lines. 
Here is the non-conical 6-arc arising from the hyperoval in $\PG(2,4).$ 
The assuciated cubic surface is the Hirschfeld-surface. 
This can be verified from the fact that the surface has 
45 Eckardt points:



\sourcecode{Hirschfeld_surface_by_arc_lifting_with_two_lines}




\bigskip


We can also create a cubic surface as an algebraic variety. 
For instance, consider the second surface  
from Dickson's list~\cite{Dickson15}.
In algebraic form, the curve is given by using the following makefile variable:

\sourcecodevariable{Dickson2_eqn_af}

The following command creates the algebraic variety defined by this equation


\sourcecode{Dickson2_group}


\bigskip


It is also possible to create a surface from an equation given as coefficient/value pairs. 
For instance, the Dickson surface $D_6$ 
can be created from the following makefile variable:


\sourcecodevariable{Dickson6}

This means that the monomials 
$m_2, m_3, m_4, m_7,$ and $m_{15}$ have a coefficient of one. 
The other monomials do not appear in the equation.
The monomial ordering is determined by the polynomial ring. 
In this case, since we use the default ring, the partition ordering is used.
The command to create the surface over $\bbF_2$ is:


\sourcecode{Dickson6_q2}

\bigskip





